\chapter{Building and Porting the Runtime}
\label{ch:runtime-build}

The earlier \emph{Building and Running} chapter (\S\ref{ch:build}) showed the
runtime from the outside: pick a profile, run \texttt{./x}, get an image. This
chapter opens the box. It explains how the three runtimes are actually
\emph{constructed}: where the Ada runtime sources come from, what makes one
profile differ from another, how the experimental \texttt{full} profile is
synthesised, and how the register packages are generated. If you ever need to
change a runtime unit, add a peripheral, or port the whole stack to another
Xtensa chip, this is the map.

\section{The idea: a board on top of \texttt{bb-runtimes}}\index{bb-runtimes}

The runtime is not written from scratch. It is grown from AdaCore's
\texttt{bb-runtimes}, the same generator that produces the bareboard GNAT
runtimes for Cortex-M, LEON, RISC-V and the rest. \texttt{bb-runtimes} knows how
to assemble a runtime from a large pool of shared Ada units; what it needs from
\emph{you} is a \emph{board definition}: the target triple, the clock, the
interrupt model, and a handful of board-specific bodies (the console, the timer,
machine reset).

\texttt{bb-runtimes} lives in the repository as a git submodule: a fork
(\texttt{rowsail/bb-runtimes}, branch \texttt{xtensa-esp32s3-port}) that carries
the ESP32-S3 board on top of the upstream sources. Keeping it a submodule means
the board port is version-pinned and the upstream history stays visible.

\cnote{There is no equivalent step in an ESP-IDF project: the C toolchain ships a
fixed newlib and FreeRTOS. Here the language runtime itself is a build product
you can read, patch and regenerate: the scheduler, the tick, the secondary
stack and the exception machinery are all Ada source in the tree.}

\section{\texttt{gen\_runtime.sh}: from board to archive}\index{gen\_runtime.sh}

One script, \texttt{crates/esp32s3\_rts/gen\_runtime.sh}, turns the board
definition into a ready-to-link runtime. It is driven by a single environment
variable:

\begin{lstlisting}[style=sh]
export ESP32S3_RTS_PROFILE=embedded      # light-tasking (default) | embedded | full
crates/esp32s3_rts/gen_runtime.sh
\end{lstlisting}

The script does four things:

\begin{enumerate}
\item \textbf{Generate the source tree.} It runs \texttt{bb-runtimes}'
  \texttt{build\_rts.py} against the \texttt{esp32s3} board, which copies the
  selected pool units plus the board's own sources into a profile-specific
  directory:
\begin{lstlisting}[style=sh]
python3 build_rts.py -f -o "$HERE" \
   --rts-src-descriptor=gnat_rts_sources/lib/gnat/rts-sources.json esp32s3
\end{lstlisting}
  The output lands in \texttt{crates/esp32s3\_rts/<profile>-esp32s3/}: a
  complete runtime with \texttt{gnat/} (sequential units) and \texttt{gnarl/}
  (tasking units) subdirectories. These directories are \emph{git-ignored}: they
  are build products, regenerated on demand.
\item \textbf{(\texttt{full} only) Apply the overlay and patches} (the next
  section).
\item \textbf{Compile} the tree with \texttt{gprbuild} against a small
  \texttt{ravenscar\_build.gpr}.
\item \textbf{Archive} the objects into \texttt{adalib/libgnat.a} (sequential)
  and \texttt{adalib/libgnarl.a} (tasking) with \texttt{xtensa-esp32-elf-ar}.
  The \texttt{esp32s3\_rts.gpr} that applications \texttt{with} points its
  \texttt{Runtime} at \texttt{<profile>-esp32s3/}, so consuming the runtime is
  just naming the crate.
\end{enumerate}

Crucially the compile-and-archive step is guarded: it runs only when
\texttt{adalib/libgnat.a} is \emph{absent}. A second invocation with the tree
already built is therefore nearly instant, but it is also the source of a trap
(\S\ref{sec:stale-cache}).

\section{Three profiles, three \texttt{system.ads}}\index{runtime profile}

What actually distinguishes the profiles is mostly one file: \texttt{system.ads},
the package that declares the target's restrictions and capabilities. The board
class maps each profile name to a variant:

\begin{center}
\begin{tabular}{l l}
\toprule
\textbf{Profile} & \textbf{\texttt{system.ads} variant} \\ \midrule
\texttt{light-tasking} & \texttt{system-xi-xtensa-light-tasking.ads} \\
\texttt{embedded}      & \texttt{system-xi-xtensa-embedded.ads} \\
\texttt{full}          & synthesised locally (see below) \\
\bottomrule
\end{tabular}
\end{center}

The \texttt{light-tasking} variant carries the restrictive pragmas:
\texttt{No\_Exception\_Propagation}, \texttt{No\_Finalization},
\texttt{No\_Secondary\_Stack}, all under \texttt{pragma Profile (Jorvik)}. The
\texttt{embedded} variant is the \emph{same Jorvik profile} with those three
restrictions lifted: exceptions propagate (zero-cost, ZCX), controlled types
finalize, and a secondary stack exists. Changing profile is, at heart, swapping
this declaration. Everything downstream (which pool units \texttt{build\_rts}
selects, which code the compiler generates) follows from it.

The board's own bodies are shared across profiles and named with the
\texttt{\_\_esp32s3} convention that \texttt{bb-runtimes} uses for board files:

\begin{center}
\begin{tabular}{l l}
\toprule
\textbf{File} & \textbf{Role} \\ \midrule
\texttt{s-bbpara\_\_esp32s3.ads} & board parameters: 240\,MHz clock, interrupt range, stacks \\
\texttt{s-bbbosu\_\_esp32s3.adb} & board support: the tick (\texttt{CCOMPARE}), interrupt dispatch \\
\texttt{s-textio\_\_esp32s3.adb} & console \texttt{Put} via the ROM \texttt{esp\_rom\_printf} \\
\texttt{s-macres\_\_esp32s3.adb} & machine reset \\
\texttt{a-intnam\_\_esp32s3.ads}  & \texttt{Ada.Interrupts.Names} for the CPU's interrupts \\
\bottomrule
\end{tabular}
\end{center}

\section{The \texttt{full} profile: overlay plus patches}\index{full profile}

There is no upstream Jorvik-free Xtensa runtime, so \texttt{full} is built on top
of \texttt{embedded} rather than generated directly. \texttt{gen\_runtime.sh}
generates the embedded tree, then layers two things on it:

\begin{itemize}
\item an \textbf{overlay} (\texttt{full\_overlay/}): a de-restricted
  \texttt{system.ads} and a set of GNARL units (interrupt management, the OS
  primitives, the tasking primitives) that replace the Jorvik-limited versions
  with the full ones; and
\item a small set of \textbf{source patches}, kept as ordinary \texttt{.patch}
  files and applied with GNU \texttt{patch}:
\end{itemize}

\begin{lstlisting}[style=sh]
patch -p1 -d "$RTS" < full_overlay/patches/01-...-secondary-stacks.patch
\end{lstlisting}

Each patch fixes one concrete defect that only surfaces under full tasking, and
each is documented at its head:

\begin{center}
\begin{tabular}{p{4.6cm} p{8.6cm}}
\toprule
\textbf{Patch} & \textbf{What it fixes} \\ \midrule
secondary-stack reclaim & free a task's heap secondary stack when it terminates \\
trampoline IRAM alias & point nested-subprogram trampolines at the executable
  (IRAM) alias of SRAM, not the data view (otherwise the first call faults) \\
\texttt{To\_Time\_Span} rounding & a nonzero \texttt{Duration} must round to
  $\geq 1$ tick (the 4.17\,ns tick is coarser than \texttt{Duration'Small}) \\
dynamic-priority re-apply & re-apply a \texttt{Set\_Priority} after the lock is
  released, which the bareboard unlock would otherwise undo \\
FIFO-within-priority & implement RM~D.2.3 so a base-priority change moves a task
  to the tail of its ready queue \\
\bottomrule
\end{tabular}
\end{center}

Keeping these as diffs rather than wholesale copied files means each change is
small, reviewable, and obviously attributable to a specific runtime behaviour,
and it keeps \texttt{gen\_runtime.sh} readable.

\section{Register packages from the SVD}\index{svd2ada}\index{SVD}

The HAL never hard-codes a peripheral address. Every register is a typed Ada
record drawn from \texttt{ESP32S3\_Registers.*}, and those packages are
\emph{generated} from Espressif's machine-readable SVD (System View Description)
by \texttt{svd2ada}. The generator script,
\texttt{libs/esp32s3\_hal/regenerate.sh}:

\begin{enumerate}
\item fetches a pinned \texttt{esp32s3.svd} (a specific \texttt{espressif/svd}
  commit, overridable with \texttt{ESP32S3\_SVD=...});
\item builds \texttt{svd2ada} \emph{from AdaCore source} (the Alire-indexed
  build is too old and mishandles dimensioned register arrays);
\item runs it into \texttt{libs/esp32s3\_hal/svd/}:
\begin{lstlisting}[style=sh]
svd2ada esp32s3.svd -o svd -p ESP32S3_Registers --boolean
\end{lstlisting}
\end{enumerate}

The root package is named \texttt{ESP32S3\_Registers} on purpose: GNAT forbids
user-defined children of \texttt{Interfaces}, so the natural
\texttt{Interfaces.ESP32S3} is not allowed. One wrinkle needs a post-process:
the SVD has a peripheral literally called \texttt{SYSTEM}, which becomes
\texttt{ESP32S3\_Registers.SYSTEM}, and because Ada is case-insensitive, a
bare \texttt{System.\dots} reference inside that unit resolves to the wrong
package. A one-line \texttt{sed} rewrites those references to the fully qualified
\texttt{Standard.System.\dots}. Add a peripheral and the recipe is the same: add
\texttt{ESP32S3.<Peri>} to the HAL, backed by the generated register record.

\section{The stale-cache trap}\index{stale cache}
\label{sec:stale-cache}

Because the heavy compile/archive step runs only when
\texttt{adalib/libgnat.a} is missing, editing a runtime source \emph{in the
generated tree} does not, on its own, rebuild anything: the next build links
the old archive and your change silently does nothing. This has burned us more
than once, so it is worth stating plainly:

\begin{quote}
\textbf{After editing a generated runtime \texttt{.adb}, delete the archives
before rebuilding:}
\begin{lstlisting}[style=sh]
rm crates/esp32s3_rts/<profile>-esp32s3/adalib/libgnat.a \
   crates/esp32s3_rts/<profile>-esp32s3/adalib/libgnarl.a
\end{lstlisting}
\end{quote}

(Editing a source under \texttt{bb-runtimes/} or \texttt{full\_overlay/} instead
regenerates correctly because the whole tree is rebuilt; the trap is specifically
editing the \emph{regenerated copy}.) The same caution applies to verifying any
runtime fix: confirm the change is in the linked archive, not just on disk,
before drawing a conclusion.

\section{Porting to another board}

The structure makes the porting checklist short. To bring the runtime to another
Xtensa SoC you would: add a board class under \texttt{bb-runtimes/xtensa/} (target
triple, clock, interrupt range, SMP flag); supply the five \texttt{\_\_<board>}
bodies (parameters, board support, console, reset, interrupt names); write the
two \texttt{system-xi-xtensa-*.ads} variants (or reuse these if the restriction
set is the same); regenerate the register packages from that chip's SVD; and
point \texttt{gen\_runtime.sh} at the new board name. The boot path (the
2nd-stage loader, the clock and cache bring-up, PSRAM) is the subject of the
next chapter.
